% !PRINTED 735
\weberpage{735}

Dass nämlich $T:T_0 = C$ eine ganze Zahl ist, folgt aus der Formel (13).

Wenn wir aber den Satz 3. auf alle Elemente des Fundamentalsystems $\mathcal{E}_i(r)$ anwenden, so ergiebt sich ein System ganzer Zahlen $e_{i,k}$ und eine ungerade Zahl $e$, so dass

\[
\mathcal{E}_i(r)^e = \pm\, \tau_0^{e_{i,0}}\tau_1^{e_{i,1}} \cdots \tau_{\nu'-1}^{e_{i,\nu'-1}}.
\]

Nehmen wir hiervon die Logarithmen, und bilden dann den Regulator $T_0$, so ergiebt sich

\[
e^{\nu'} T_0 = \pm\, \Sigma \pm e_{0,0}\, e_{1,1} \cdots e_{\nu'-1,\nu'-1}\, T,
\]

folglich

\[
e^{\nu'} = \pm\, C\, \Sigma \pm e_{0,0}\, e_{1,1} \cdots e_{\nu'-1,\nu'-1},
\]

und da hierin die linke Seite ungerade ist, so kann $C$ nicht gerade sein, w. z. b. w.

Wenn alle conjugirten Werthe der Normaleinheit $\mathcal{E}(r_\beta)$ einerlei Zeichen haben, so gilt dasselbe auch von $\mathcal{E}(r_\beta)^e$, und folglich müssen in diesem Falle in der Formel (20) die Exponenten $e_0, e_1, \ldots, e_{\nu'-1}$ nach 2. gerade Zahlen sein. Schreibt man dann die Formel (20) so:

\[
\mathcal{E}(r) = \pm\, \tau_0^{e_0}\tau_1^{e_1} \cdots \tau_{\nu'-1}^{e_{\nu'-1}}\, \mathcal{E}(r)^{1-e},
\]

so erhält man daraus den Satz:

5. \gesperrt{Eine Normaleinheit des Körpers $H_m$, die mit ihren Conjugirten gleiches Vorzeichen hat, ist, vom Vorzeichen abgesehen, ein Quadrat einer Normaleinheit.}

\begin{center}
\S.\,201.

\gesperrt{Fundamentalsystem von Einheiten des Körpers $H_m$.}
\end{center}

Es bleibt noch übrig, den Nenner $D$ in der Formel für $B$ [\S.\,197, (6)] zu bestimmen, der durch die Gleichung

\[
E = E'\, D
\]

definirt war, worin $E$ und $E'$ die Regulatoren der Körper $H_m, H_\mu$ bedeuteten.

Hier handelt es sich also nicht mehr um die Normaleinheiten, sondern um die Einheiten des Körpers $H_m$ überhaupt.
